% !TEX TS-program = pdflatex
% !TeX program = pdflatex
% !TEX encoding = UTF-8
% !TEX spellcheck = en_US

\documentclass[xcolor=table]{beamer}

\usepackage{../extra/beamer/karimnlp}

%\usepackage{fullpage}
%\usepackage[left=2.8cm,right=2.2cm,top=2 cm,bottom=2 cm]{geometry}
\setbeamersize{text margin left=10pt,text margin right=10pt}



%\usepackage[T1]{fontenc}
\usepackage{textcomp}
\usepackage[utf8]{inputenc}
%\usepackage[french,english]{babel}
\usepackage{arabtex}
\usepackage{txfonts}
\usepackage{tipa}
\usepackage[]{graphicx}
\usepackage{multirow}
\usepackage{hyperref}
\usepackage{colortbl}
\usepackage{tabularray}
\usepackage{listingsutf8}
\usepackage{wrapfig}
\usepackage{multicol}
\usepackage[export]{adjustbox} %for images in table, also for frame
\usepackage[many]{tcolorbox}
\usepackage{wasysym}
\usepackage[lined]{algorithm2e}
\usepackage{alltt} %verbatim with commands
\usepackage{longtable}
\usepackage{tabu}

\usepackage{CJKutf8}

\usepackage{amsmath,amssymb, amsfonts} 

%\usepackage{newtxtext,newtxmath}


%\usepackage{acolor}


\definecolor{lightblue}{HTML}{D0D2FF}
\definecolor{lightyellow}{HTML}{FFFFAA}
\definecolor{darkblue}{HTML}{0000BB}
\definecolor{olivegreen}{HTML}{006600}
\definecolor{darkgreen}{HTML}{008B45} %009B55
\definecolor{violet}{HTML}{6600CC}
\definecolor{deeppink}{HTML}{FF1493}
\definecolor{orangey}{HTML}{FFBB00}


\makeatletter
\def\beamer@calltheme#1#2#3{%
	\def\beamer@themelist{#2}
	\@for\beamer@themename:=\beamer@themelist\do
	{\IfFileExists{\beamer@themelocation/#3\beamer@themename.sty}{%
	\usepackage[{#1}]{\beamer@themelocation/#3\beamer@themename}}{\usepackage[{#1}]{#3\beamer@themename}}
		
	}%
}

\def\usefolder#1{
	\def\beamer@themelocation{#1}
}
\def\beamer@themelocation{}
\makeatother

\usefolder{../../extra/beamer-themes}

\usetheme{Karim} % Antibes Boadilla Warsaw Karim

\beamertemplatenavigationsymbolsempty

\let\oldcite\cite
\renewcommand{\cite}[1]{{\bfseries\color{blue}\oldcite{#1}}}

\tcbuselibrary{listings}


%\renewcommand{\baselinestretch}{1.5}

\def\supit#1{\raisebox{0.8ex}{\small\it #1}\hspace{0.05em}}

\AtBeginSection{%
	\begin{frame}
		\sectionpage
	\end{frame}
}

\newcommand{\rottext}[2]{%
	\rotatebox{90}{%
	\begin{minipage}{#1}%
		\raggedleft#2%
	\end{minipage}%
	}%
}


\institute{ %
Laboratoire de la Communication dans les Systèmes Informatiques (LCSI)
\par
École  nationale Supérieure d'Informatique (ESI), Algiers, Algeria
}
\author[Abdelkrime Aries (2023/2024)] %
{Abdelkrime Aries}
%\titlegraphic{\includegraphics[height=1cm]{../img/esi-logo.png}%\hspace*{4.75cm}~


\date{Academic year: 2023/2024} %\today

\titlegraphic{%
	\includegraphics[height=1cm]{../../img/misc/esi-logo.png}%
	\hspace{2cm}%
	\includegraphics[height=1cm]{../../img/misc/lcsi-logo.png}%
	\hspace{2cm}%
	\includegraphics[height=1.5cm]{../../img/misc/esi.nlp-logo.png}%
}

%\setbeamertemplate{headline}{}

\newcommand{\kurl}[1]{{\scriptsize\bfseries\color{orangey}\url{#1}}}

\newcommand{\keyword}[1]{\textcolor{red}{\bfseries\itshape #1}}
\newcommand{\expword}[1]{\textcolor{olivegreen}{#1}}
\newcommand{\optword}[1]{\textcolor{violet}{\bfseries #1}}

\makeatletter
\newcommand\mysphere{%
	\parbox[t]{10pt}{\raisebox{0.2pt}{\beamer@usesphere{item projected}{bigsphere}}}}
\makeatother

%\let\oldtabular\tabular
%\let\endoldtabular\endtabular
%\renewenvironment{tabular}{\rowcolors{2}{white}{lightblue}\oldtabular\rowcolor{blue}}{\endoldtabular}


%\NoAutoSpacing %french autospacing after ":"

\def\graphpath{}

\newcommand{\changegraphpath}[1]{\def\graphpath{#1}}


\newcommand{\vgraphpage}[2][.84\textheight]{%
%	\begin{center}%
		\includegraphics[height=#1]{\graphpath #2}%
%	\end{center}%
}

\newcommand{\hgraphpage}[2][\textwidth]{%
%	\begin{center}%
		\includegraphics[width=#1]{\graphpath #2}%
%	\end{center}%
}

\newcommand{\graphpage}[2][]{%
	\includegraphics[#1]{\graphpath #2}%
}

\bibliographystyle{apalike}

\newcommand{\insertbibliography}[2]{
	\appendix
	\section*{Bibliography}
	\nocite{#2}
%	\makeatletter % to change template
%	\setbeamertemplate{headline}[default] % not mandatory, but I though it was better to set it blank
%	\def\beamer@entrycode{\vspace*{-\headheight}} % here is the part we are interested in :)
%	\makeatother
	\begin{multicols*}{2}[\frametitle{\insertsection} \usebeamertemplate*{frametitle}]%\usebeamertemplate*{frametitle}\frametitle{Références}
		\tiny
		\bibliography{#1}
	\end{multicols*}
}

\definecolor{my-grey}{RGB}{233, 233, 233}

\newcommand{\insertlicence}{
	\begin{frame}[plain]
	\frametitle{License: CC-BY 4.0}

	\begin{tcolorbox}[colback=cyan,
		colframe=cyan,  
		arc=0pt,outer arc=0pt,
		valign=top, 
		halign=center,
		width=\textwidth]
		
		\includegraphics[width=.5cm]{../../img/licence/cc_icon_white_x2.png}
		\includegraphics[width=.5cm]{../../img/licence/attribution_icon_white_x2.png}
		
		\color{white}
		\bfseries Attribution 4.0 International (CC BY 4.0) \\
		\tiny \url{https://creativecommons.org/licenses/by/4.0/deed.en}
		
	\end{tcolorbox}\vspace{-.5cm}
	\begin{tcolorbox}[colback=my-grey,
		colframe=my-grey,  
		center, arc=0pt,outer arc=0pt,
		valign=top, 
		halign=left,
		width=\textwidth]
		
		\tiny
		
		\begin{center}
			\bfseries\Large
			You are free to:
		\end{center}
		
		\begin{minipage}{0.83\textwidth}
			\begin{itemize}
				\item[] \textbf{Share} — copy and redistribute the material in any medium or format
				\item[] \textbf{Adapt} — remix, transform, and build upon the material
				for any purpose, even commercially.
			\end{itemize}
		
			\vspace{12pt}
		
			\textit{The licensor cannot revoke these freedoms as long as you follow the license terms.}
		\end{minipage}
		\begin{minipage}{0.15\textwidth}
			\includegraphics[width=\textwidth]{../../img/licence/FreeCulturalWorks_seal_x2.jpg}
		\end{minipage}
	
		
		\begin{center}
			\bfseries\Large
			Under the following terms:
		\end{center}
		
		\begin{itemize}
			\item[] \textbf{Attribution} — You must give appropriate credit, provide a link to the license, and indicate if changes were made. You may do so in any reasonable manner, but not in any way that suggests the licensor endorses you or your use.
			\item[] \textbf{No additional restrictions} — You may not apply legal terms or technological measures that legally restrict others from doing anything the license permits.
		\end{itemize}
		
	\end{tcolorbox}
	
%	\begin{center}
%		\bfseries Attribution 4.0 International (CC BY 4.0)
%		\url{https://creativecommons.org/licenses/by/4.0/deed.fr}
%	\end{center}

%	\tiny
%
%	Vous êtes autorisé à : 
%	\begin{itemize}
%		\item \textbf{Partager} — copier, distribuer et communiquer le matériel par tous moyens et sous tous formats
%		\item \textbf{Adapter} — remixer, transformer et créer à partir du matériel
%	\end{itemize}
%	
%	Selon les conditions suivantes : 
%	\begin{itemize}
%		\item \textbf{Attribution} — Vous devez créditer l'Œuvre, intégrer un lien vers la licence et indiquer si des modifications ont été effectuées à l'Oeuvre. Vous devez indiquer ces informations par tous les moyens raisonnables, sans toutefois suggérer que l'Offrant vous soutient ou soutient la façon dont vous avez utilisé son Oeuvre.
%		\item \textbf{Pas d'Utilisation Commerciale} — Vous n'êtes pas autorisé à faire un usage commercial de cette Oeuvre, tout ou partie du matériel la composant. 
%		\item \textbf{Pas de restrictions complémentaires} — Vous n'êtes pas autorisé à appliquer des conditions légales ou des mesures techniques qui restreindraient légalement autrui à utiliser l'Oeuvre dans les conditions décrites par la licence.
%	\end{itemize}

	\end{frame}
}

\settowidth{\leftmargini}{\usebeamertemplate{itemize item}}
\addtolength{\leftmargini}{\labelsep}

\AtBeginDocument{
	\newcolumntype{L}[2]{>{\vbox to #2\bgroup\vfill\flushleft}p{#1}<{\egroup}} 
	
	\begin{frame}[plain]
		\maketitle
	\end{frame}

	\insertlicence
}


% needs etoolbox; to break links after -
\appto\UrlBreaks{\do\-}


\makeatletter
\newcommand{\xRightarrow}[2][]{\ext@arrow 0359\Rightarrowfill@{#1}{#2}}
\makeatother


\usefonttheme{structurebold}
%\usefonttheme{professionalfonts}


% figure caption
\setlength\abovecaptionskip{2pt}

% equations
\AtBeginDocument{
\abovedisplayskip=0pt
\abovedisplayshortskip=0pt
\belowdisplayskip=0pt
\belowdisplayshortskip=0pt
}


\subtitle[09- Coreference resolution]{Chapter 09\\Coreference resolution} 

\changegraphpath{../img/coref/}

\begin{document}
	
\begin{frame}
	\frametitle{\inserttitle}
	\framesubtitle{\insertshortsubtitle: Introduction}

	\begin{exampleblock}{Example of a paragraph}
		\begin{center}
			\Large\bfseries
		The cat saw a flower. \textbf{\color{red}It} smelled \textbf{\color{red}it}. \textbf{\color{red}It} has a very nice scent.
		\end{center}
	\end{exampleblock}
	
	\begin{itemize}
		\item Who smelled whom: the cat or the flower?
		\item Who has a nice scent: the cat or the flower?
	\end{itemize}

\end{frame}


\begin{frame}
	\frametitle{\inserttitle}
	\framesubtitle{\insertshortsubtitle: Plan}

	\begin{multicols}{2}
	%	\small
	\tableofcontents
	\end{multicols}

\end{frame}

%===================================================================================
\section{References}
%===================================================================================

\begin{frame}
	\frametitle{\insertshortsubtitle}
	\framesubtitle{\insertsection}
	
\end{frame}

\subsection{References forms}

\begin{frame}
	\frametitle{\insertshortsubtitle: \insertsection}
	\framesubtitle{\insertsubsection}
	
	\begin{itemize}
		\item \optword{Pronouns}
		\begin{itemize}
			\item Personal: \expword{\underline{Karim} entered. \underline{He} started class.}
			\item Possessive: \expword{\underline{Karim} has started \underline{his} course.}
			\item ...
		\end{itemize}
	
		\item \optword{Nominal phrases}
		\begin{itemize}
			\item \expword{I have a little \underline{cat}. \underline{This animal} is very mean.}
		\end{itemize}
	
		\item \optword{Proper nouns} 
		\begin{itemize}
			\item \expword{The \underline{école nationale supérieure d'informatique} is located in Algiers. Like all Algerian universities, you must have the baccalaureate to study at \underline{ESI}.}
		\end{itemize}
	
		\item \optword{Zero anaphora}
		\begin{itemize}
			\item \expword{\underline{Karim} presented and \underline{$ \phi $} explained the course.}
			\item \expword{\begin{CJK}{UTF8}{min}\underline{カリムさん}はESIに生きます。\underline{$ \phi $} あそこに教えています。\end{CJK}}
		\end{itemize}
	\end{itemize}
	
\end{frame}

\subsection{Referencing manner}

\begin{frame}
	\frametitle{\insertshortsubtitle: \insertsection}
	\framesubtitle{\insertsubsection}
	
	\begin{itemize}
		\item \optword{Anaphora}
		\begin{itemize}
			\item a reference to a previous word/phrase is called \keyword{antecedent}.
			\item E.g. \expword{\underline{The lesson} is very long. \underline{It} will take longer.}
		\end{itemize}
	
		\item \optword{Cataphora}
		\begin{itemize}
			\item a reference to a following word/phrase is called \keyword{postcedent}.
			\item E.g. \expword{\underline{It} is very long, \underline{this course}!}
		\end{itemize}
	
		\item \optword{Shared antecedents}
		\begin{itemize}
			\item a reference to several words and/or phrases.
			\item E.g. \expword{\underline{The lesson} will be followed by \underline{an exercise}. \underline{They} are important to understand the subject.}
		\end{itemize}
	
		\item \optword{Coreference noun phrases}
		\begin{itemize}
			\item two noun phrases where each is a reference to the other.
			\item E.g. \expword{\underline{Some of my colleagues} have been really supportive. \underline{This kind of people} earns my gratitude.}
		\end{itemize}
	\end{itemize}
	
\end{frame}

\subsection{Properties of Coreference Relations}

\begin{frame}
	\frametitle{\insertshortsubtitle: \insertsection}
	\framesubtitle{\insertsubsection\ (1)}
	
	
	\begin{itemize}
		\item \optword{Number}: Singular, Dual, Plural.
		\begin{itemize}
			\item \expword{\underline{Lions} are chasing a \underline{dear}. \underline{They} are so fast.}
		\end{itemize}
	
		\item \optword{Person}: First, Second, Third.
		\begin{itemize}
			\item \expword{\underline{My brother}\textsubscript{1} and \underline{I}\textsubscript{2} have repaired \underline{his bicycle}\textsubscript{1} after \underline{mine}\textsubscript{2}.}
		\end{itemize}
		
		\item \optword{Gender}: Masculine, Feminine, Neutral.
		\begin{itemize}
			\item \expword{When the girl met her \underline{father}, \underline{he} was very happy.}
		\end{itemize}
		
		\item \optword{Binding theory's constraints}: syntactic constraints on the mention-antecedent relation.
		\begin{itemize}
			\item \expword{Janet bought herself a bottle of fish sauce.} [herself = Janet]
			\item \expword{Janet bought her a bottle of fish sauce.} [her $\ne$ Janet]
		\end{itemize}
	
	\end{itemize}
	
\end{frame}

\begin{frame}
	\frametitle{\insertshortsubtitle: \insertsection}
	\framesubtitle{\insertsubsection\ (2)}
	
	\begin{itemize}
		
		\item \optword{Recency}: entities mentioned more recently are more likely to be an antecedent.
		\begin{itemize}
			\item \expword{The doctor found an old map. Jim found an even older \underline{map}. \underline{It} is about an island.}
		\end{itemize}
		
		\item \optword{Grammatical relations}: subjects are more likely to be an antecedents than objects.
		\begin{itemize}
			\item \expword{\underline{Karim} went to a restaurant with his friend. \underline{He} asked for a couscous dish.}
		\end{itemize}
		
		\item \optword{Verb semantics}: the mention follows the emphasis of the verb (the one that caused the event).
		\begin{itemize}
			\item \expword{\underline{John} telephoned Bill. \underline{He} lost the laptop.}
			\item \expword{John criticized \underline{Bill}. \underline{He} lost the laptop.}
		\end{itemize}
	\end{itemize}
	
\end{frame}

\begin{frame}
	\frametitle{\insertshortsubtitle: \insertsection}
	\framesubtitle{\insertsubsection: Some humor}
	
	\begin{center}
		\vgraphpage{humor/humor-ref.jpg}
	\end{center}
	
\end{frame}

%===================================================================================
\section{Coreference resolution}
%===================================================================================

\begin{frame}
	\frametitle{\insertshortsubtitle}
	\framesubtitle{\insertsection}
	
	\hgraphpage{coref-arch.pdf}
	
	\begin{itemize}
		\item \optword{Mention detection}
		\item \optword{Linking models}
		\begin{itemize}
			\item Mention-Pair: each mention with others.
			\item Mention-Rank: a mention must be linked to only one antecedent.
			\item Entity-based: detect coreference clusters.
		\end{itemize}
		
		\item \optword{Coreference resolution approaches}
		\begin{itemize}
			\item Rules-based
			\item Features-based (manually engineered)
			\item Embedding-based
		\end{itemize}
	\end{itemize}
	
\end{frame}


\subsection{Mention detection}

\begin{frame}
	\frametitle{\insertshortsubtitle: \insertsection}
	\framesubtitle{\insertsubsection}
	
	\begin{itemize}
		\item \optword{Entity extraction}
		\begin{itemize}
			\item Nominal phrases, pronouns, named entities.
		\end{itemize}
	
		\item \optword{Entities filtering}
		\begin{itemize}
			\item Using linguistic rules. E.g. In the expression ``\expword{It is thought that ...}", the pronoun ``\expword{It}" is not a reference. In this case, a list of cognitive verbs \expword{believe, think, etc.} can be used to filter this pronoun.
			\item Using machine learning: one detector for references and another for antecedents.
		\end{itemize}
	
	\end{itemize}
	
\end{frame}


\begin{frame}
	\frametitle{\insertshortsubtitle: \insertsection}
	\framesubtitle{\insertsubsection: Rule-based \cite{2013-lee-al}}
	
	\begin{itemize}
		\item Delete a mention if there is another larger one with the same syntactic head. 
		E.g., ``\expword{the software company}" vs. ``\expword{the software company situated in Algiers}". 
		
		\item Discard numerical entities such as percentages, money, cardinals and quantities.
		E.g., ``\expword{9\%}", ``\expword{100\$}", ``\expword{one hundred nine}".
		
		\item Eliminate mentions that are part of another (quantifier + ``of"). 
		E.g., ``\expword{a total of 177 projects}", ``\expword{none of them}", ``\expword{millions of people}".
		
		\item Remove pronouns that are not used as a reference (pleonastic ``it"). 
		E.g., ``\expword{It is thought that ...}", the pronoun ``it" is not a reference. 
		A list of cognitive verbs \expword{believe, think, etc.} can be used to filter this pronoun.
		
		\item Eliminate adjectival forms of nations or national acronyms.
		E.g., ``\expword{American}", ``\expword{U.S.}", ``\expword{U.K.}".
		
	\end{itemize}
	
\end{frame}


\begin{frame}
	\frametitle{\insertshortsubtitle: \insertsection}
	\framesubtitle{\insertsubsection: NN-based \cite{2020-yu-al}}
	
	\begin{figure}[ht]
		\centering
		\hgraphpage[\textwidth]{mention-detection-arch.pdf}
		\caption[Three architectures to detecting mentions]{Three architectures to detecting mentions; Inspired from \cite{2020-yu-al}}
		\label{fig:det-mention-yu}
	\end{figure}
	
\end{frame}



\subsection{Linking}

\begin{frame}
	\frametitle{\insertshortsubtitle: \insertsection}
	\framesubtitle{\insertsubsection: Mention-Pair model}
	
	\begin{itemize}
		\item Detect if there is a coreference between two mentions.
		\item Binary classification (coreference, not coreference).
		\item \textbf{Problem}: contradictions. E.g. \expword{Ms Kennedy $ \leftarrow $ Kennedy, Kennedy $ \leftarrow $ He}
	\end{itemize}
	\begin{figure}
		\centering
		\hgraphpage[.8\textwidth]{mention-pair-exp.pdf}
		\caption{Example of Mention-Pair model \cite{2019-jurafsky-martin}}
	\end{figure}
	
\end{frame}

\begin{frame}
	\frametitle{\insertshortsubtitle: \insertsection}
	\framesubtitle{\insertsubsection: Mention-Rank model}
	
	\begin{itemize}
		\item Among previous mentions, decide which one is the antecedent.
		\item $ \hat{a} = \arg\max\limits_{a_i \in \{\epsilon, a_1, \ldots, a_{n-1}\}} P(a_i|a_n) $
		\item $ \epsilon $: means, there is no antecedent.
	\end{itemize}
	\begin{figure}
		\centering
		\hgraphpage{mention-rank-exp.pdf}
		\caption{Example of Mention-Rank model \cite{2019-jurafsky-martin}}
	\end{figure}
	
\end{frame}

\begin{frame}
	\frametitle{\insertshortsubtitle: \insertsection}
	\framesubtitle{\insertsubsection: Entity-based model}
	
	\begin{itemize}
		\item Also called: cluster-ranking model
		\item \textbf{Input}: two mentions' clusters
		\item verify if the two clusters are compatible
		\item Estimate if a cluster is antecedent to another such as in Mention-Rank
		\item If the two clusters represent the same mention, merge them
		\item \textbf{Output}: A set of clusters
	\end{itemize}
	
\end{frame}

\subsection{Evaluation}

\begin{frame}
	\frametitle{\insertshortsubtitle: \insertsection}
	\framesubtitle{\insertsubsection}
	
	\begin{itemize}
		\item \keyword{MUC}: \optword{Message Understanding Conference}
		\begin{itemize}
			\item Link-based metric
			\item Number of shared binary links between the reference and the system
		\end{itemize}
		\item \keyword{B\textsuperscript{3}}
		\begin{itemize}
			\item Mention-based metric
			\item Global R/P is calculated based on R/P of individual mentions.
		\end{itemize}
		\item \keyword{CEAF}: \optword{Constrained entity-alignment F-Measure}
		\begin{itemize}
			\item mention-based OR entity-based (2 versions depending on used similarity)
			\item A mention of the system is aligned with a single reference
		\end{itemize}
		\item \keyword{BLANC} 
		\begin{itemize}
			\item Link-based metric
			\item Global R/P is calculated based on R/P of coreference links and that of non coreference
		\end{itemize}
		\item \keyword{LEA}: \optword{Link based entity aware}
		\begin{itemize}
			\item Entity size as an importance measure
			\item Resolved coreferences' links are evaluated.
		\end{itemize}
	\end{itemize}
	
\end{frame}

\begin{frame}
	\frametitle{\insertshortsubtitle: \insertsection}
	\framesubtitle{\insertsubsection: Some humor}
	
	\begin{center}
		\vgraphpage{humor/humor-res.jpg}
	\end{center}
	
\end{frame}

%===================================================================================
\section{Related tasks}
%===================================================================================

\begin{frame}
	\frametitle{\insertshortsubtitle}
	\framesubtitle{\insertsection}
	
	\begin{itemize}
		\item Used tasks
		\begin{itemize}
			\item Preprocessing task (Chapter 3)
			\item PoS tagging (Chapter 5)
			\item Named entities recognition (Chapter 5)
		\end{itemize}
		\item Similar tasks
		\begin{itemize}
			\item Entity linking
			\item Citation attribution: find who said/wrote a discourse
		\end{itemize}
	\end{itemize}
	
\end{frame}

\subsection{Entity linking}

\begin{frame}
	\frametitle{\insertshortsubtitle: \insertsection}
	\framesubtitle{\insertsubsection}
	
	\begin{itemize}
		\item Associate a mention in a text with its representation in a structured knowledge base.
		\item useful to have in-depth knowledge on some entities in the real world.
		\item \keyword{Wikification}: link a mention with a Wikipédia page
		\item Entity linking steps: 
		\begin{itemize}
			\item \optword{Mention detection}: detect all knowledge base entities linked to a mention using queries.
			\item \optword{Mention disambiguation}: find the most likely entity using machine learning.
		\end{itemize}
		\item Some APIs: 
		\begin{itemize}
			\item \url{https://labs.tib.eu/falcon/}
			\item \url{https://dandelion.eu}
		\end{itemize}
	\end{itemize}
	
\end{frame}

\subsection{Named entity recognition (NER)}

\begin{frame}
	\frametitle{\insertshortsubtitle: \insertsection}
	\framesubtitle{\insertsubsection}
	
	\begin{itemize}
		\item locate and classify named entities in a text
		\item \keyword{Named entity}: persons, places, organizations, quantities, etc.
		\item a data extraction subtask
		\item Recognition techniques
		\begin{itemize}
			\item \optword{Rules}: use rules and lists of names to search for entities and detect their types.
			\item \optword{Features-based ML}: word embeddings, prefixes and suffixes, list membership, neighbors features, etc.
			\item \optword{Sequences labeling}: As seen in chapter 5 using \keyword{IOB} annotation.
			
			\expword{\scriptsize
			$ \underbrace{Google}_{B-ORG} $ 
			$ \underbrace{LLC}_{I-ORG} $ 
			$ \underbrace{is}_{O} $ 
			$ \underbrace{founded}_{O} $ 
			$ \underbrace{in}_{O} $ 
			$ \underbrace{Silicon}_{B-LOC} $ 
			$ \underbrace{Valley}_{B-LOC} $ 
			$ \underbrace{by}_{O} $ 
			$ \underbrace{Larry}_{B-PER} $ 
			$ \underbrace{Page}_{I-PER} $ 
			$ \underbrace{and}_{O} $ 
			$ \underbrace{Sergey}_{B-PER} $ 
			$ \underbrace{Brin}_{I-PER} $
			}
		\end{itemize}
	\end{itemize}
	
\end{frame}


\insertbibliography{NLP09}{*}

\end{document}

